\small\noindent\textbf{초록}\quad
본 연구는 이상적인 완전혼합 수용액 열 공정에서 PID, discrete linear-quadratic regulator(LQR), linear-quadratic integral(LQI) control을 비교한다. Rust로 구현한 정확한 zero-order hold(ZOH) 전파, 독립적으로 확인한 Riccati 해, equal-budget tuning, $\pm30\%$ 불확실성, constant thermal load, sensor noise, actuator saturation, sensor lag와 measurement delay를 평가했다. Equal-budget tuning은 공칭 integral absolute error(IAE) 순위를 LQI--PID--LQR에서 LQR--LQI--PID로 바꾸지만, 튜닝된 LQR은 \SI{20}{s} sensor lag와 \SI{10}{s} measurement delay에서 saturation의 지배를 크게 받는다. Integral action은 실현 가능한 constant load의 offset을 제거한다. 정상상태 noise 차이는 통계적으로 식별되지만 실용적으로 작다. Validated interval bound로 비포화 LQR의 안정성을 인증하고, 세분된 interval Jury criterion으로 선언한 연속 parameter box 전체에서 LQI 안정성을 인증했다. PID는 formal interval proof 대신 dense-grid evidence를 갖는다. 본 연구의 기여는 조건부 성능 trade-off의 재현 가능한 규명이며, 새로운 제어기나 물리 실험 또는 보편적 강인성 주장이 아니다.

\par\vspace{0.7em}\noindent\textbf{Keywords}\quad PID control; linear-quadratic control; integral action; parametric uncertainty; thermal process; computational reproducibility.
